% Nag is used for debug
\RequirePackage[l2tabu, orthodox]{nag}
\documentclass[a4paper,french,twoside,10pt]{report}
%\documentclass[a4paper,twoside,10pt]{article}
%\usepackage[bookman]{quotchap}

\usepackage[bindingoffset=1.5cm,centering,includeheadfoot,margin=3cm]{geometry}

%\frenchspacing
\usepackage[frenchb]{babel}
\usepackage[T1]{fontenc}
\usepackage[utf8]{inputenc}
\usepackage[pdftex]{graphicx}
%\usepackage{subfig}

\usepackage[protrusion=true,expansion=true]{microtype}

\usepackage{fancybox}
\usepackage{shadow}
\usepackage{xspace}
\usepackage{acronym}

% Maths enhancing
\usepackage{amsmath}
\usepackage{amssymb}
\usepackage{amsthm}
\usepackage{dsfont}

\usepackage{algorithm2e}

% JP : automatic line return after paragraph
\makeatletter
\renewcommand\paragraph{\@startsection{paragraph}{4}{\z@}%
    {-3.25ex\@plus -1ex \@minus -.2ex}%
    {1.5ex \@plus .2ex}%
{\normalfont\normalsize\bfseries}}
\makeatother

% JP : use amsthm package and define some useful commands
\newtheorem{theorem}{Theorem}
\newtheorem{definition}{Definition}
\newtheorem{proposition}{Proposition}

%\usepackage[right]{eurosym}
\usepackage{lmodern} % new latin extended computer modern font}, use with T1
%\usepackage{pxfonts} % has replaced palatino
%\usepackage{cmbright}

\usepackage{listings}

%\usepackage{booktabs}
%\usepackage{array}

\usepackage[allowlitunits]{siunitx}
% Try $20\micro\seconds$

%\addtolength{\oddsidemargin}{-1cm}
%\addtolength{\evensidemargin}{-1cm}
%\addtolength{\textwidth}{+2cm}
%\usepackage{showframe}

%\usepackage{lastpage}
% for \pageref{LastPage}

% fancy header
\usepackage{fancyhdr}
\pagestyle{fancy}
%\pagestyle{plain} % Imprime simplement le numero de page centre en bas
%\pagestyle{headings}
\fancyhf{} % clear default header and footer

\fancyhead[LE,RO]{\textit{\nouppercase{\leftmark}}}
%\fancyhead[LO,RE]{\textit{\nouppercase{\thepage{} of \pageref{LastPage}}}}
\fancyhead[CE,CO]{}
\fancyfoot[C]{--\thepage--}
\setlength{\headheight}{14pt}

\newcommand\Hrule{\noindent\rule{\linewidth}{1.5pt}}
%\AL: to add abs and norm for fracs
%\DeclarePairedDelimiter\abs{\lvert}{\rvert}%
%\DeclarePairedDelimiter\norm{\lVert}{\rVert}%
\usepackage[pdftex,
    pdfauthor={Alexandre Lafaye},
    pdftitle={Laboratoire MUI},
    pdfsubject={},
    pdfkeywords={rapport,laboratoire,HEIG-VD,MUI,Microinformatique},
pdfcreator=pdflatex]{hyperref}
\hypersetup{colorlinks, citecolor=black, filecolor=black, linkcolor=black, urlcolor=black}

\newcommand{\weblink}[2]{\href{#1}{\url{#1}: \textsc{#2}}}

%Maths shortcuts
%\newcommand{\R}{\mathbb{R}}
%\newcommand{\E}{\mathbb{E}}
%\newcommand{\N}{\mathbb{N}}
%\newcommand{\Var}{\mathrm{Var}}
%\newcommand{\Cov}{\mathrm{Cov}}
%\newcommand{\Prob}{\mathbb{P}}
%\newcommand{\trace}{\mathbb{\text{trace}}}
%\def\mathbi#1{\textbf{\em #1}}
%%\def\mathi#1{\em #1}
%\newcommand{\Z}{\mathbi{Z}_2}
%\newcommand{\kMeans}{{\em k}-Means\xspace}
%\newcommand{\kmeans}{{\em k}-means\xspace}
%\newcommand{\kmeanspp}{{\em k}-means$++$\xspace}
%\newcommand{\captionwidth}{0.9\textwidth}

% index stuff
\usepackage{tocbibind} % Index in TOC
%\usepackage{makeidx}
%\usepackage{showidx} % show index entries in margin
%\makeindex

%\newcommand{\bib}[4]{\item{\textsc{#1}, \emph{#2}, #3\\\hspace{2em}#4}}

% ---------------------------------------------------------

\newcommand{\titleinfo}{Laboratoire MUI}
\author{Alexandre Lafaye}
\date{12 Mars 2016}
\title{\titleinfo}

% Margins for handwritten notes
%\addtolength{\oddsidemargin}{-2cm}
%\addtolength{\evensidemargin}{-2cm}
%\addtolength{\textwidth}{-2cm}

\usepackage{color}
\definecolor{Brown}{cmyk}{0,0.81,1,0.60}
\definecolor{OliveGreen}{cmyk}{0.64,0,0.95,0.40}
\definecolor{CadetBlue}{cmyk}{0.62,0.57,0.23,0}
\definecolor{lightlightgray}{cmyk}{0,0,0,0.02}
\definecolor{gray}{cmyk}{0,0,0,0.8}

% LaTeX magic: make sections and parts have a cleardoublepage
\let\stdsection\section
\renewcommand*\section{\cleardoublepage\stdsection}

\begin{document}
% JP
%\frontmatter

\lstset{
    language=Python,                             % Code langugage
    basicstyle=\ttfamily,                   % Code font, Examples: \footnotesize, \ttfamily
    keywordstyle=\color{OliveGreen},        % Keywords font ('*' = uppercase)
    commentstyle=\color{gray},              % Comments font
    numbers=left,                           % Line nums position
    numberstyle=\tiny,                      % Line-numbers fonts
    stepnumber=1,                           % Step between two line-numbers
    numbersep=5pt,                          % How far are line-numbers from code
    backgroundcolor=\color{lightlightgray}, % Choose background color
    frame=none,                             % A frame around the code
    tabsize=2,                              % Default tab size
    captionpos=b,                           % Caption-position = bottom
    breaklines=true,                        % Automatic line breaking?
    breakatwhitespace=false,                % Automatic breaks only at whitespace?
    showspaces=false,                       % Dont make spaces visible
    showtabs=false,                         % Dont make tabls visible
    morekeywords={uint32_t,uint8_t,uint16_t,time_spec_t,size_t,
    clock_config_t},
}
\pagestyle{empty}
\pagenumbering{roman}
\begin{titlepage}
    \null\vfil
    \begin{flushleft}
	\begin{center}
	    \Huge \textbf{\titleinfo}
	\end{center}
    \end{flushleft}
    \par
    \hrule height 4pt
    \par
    \begin{flushright}
	\large
	\textsl{Rapport de laboratoire} \par
    \end{flushright}
    \vspace{\fill}
    \vspace*{\stretch{1}}

    \begin{center}
	\Large
	\textsc{Alexandre Lafaye\\
	    %\normalsize
	HEIG-VD -- 2016} \\
	\href{mailto:alexandre.lafaye@epfl.ch}{alexandre.lafaye@epfl.ch}
    \end{center}
    \vspace{\fill}
    \vspace*{\stretch{2}}
    \begin{center}
	\large
	Supervised by: \\
	\vspace{\fill}
	Cedric Bornand -- HEIG-VD\\
	\href{http://heig-vd.ch/}{http://heig-vd.ch/} \\
    \end{center}

    \vspace{\fill}
    \vspace*{\stretch{2}}

    \begin{figure}[!h]
	\centering
	\parbox{1.4in}{
	\includegraphics[height=4em]{logo_heig_h80.jpg}}
    \end{figure}

    \vspace*{1cm}
\end{titlepage}

\pagestyle{fancy}

\part{Introduction}
\section*{Introduction}

\subsection*{Introduction}
Ce labo a pour but de tester les possibilites offertes par les interruptions dans le cadre de la conception de programmes pour micro-controleurs.
Pour ce travail pratique, une MSP430f5529 a été utilisée, le test est le debug effectué au travers de son interface JTAG.

\subsection{Buts}
Ce laboratoire est séparé en quatre partie:
\begin{itemize}
\item Tout d'abord changer l'état de la LED PAD5 toutes les demi-secondes
\item Ensuite, tout en laissant la LED clignoter, faire changer l'état des LEDs PAD[1,2,3] avec le bouton S1 ou S2
\item Recopier l'état du bouton S1 sur les LEDs PAD[1,2,3]
\item Finalement faire changer l'état des LEDs PAD[1,2,3] en appuyant sur le bouton S1
\end{itemize}


\part{Conception}
\section{Conception}
\subsection{Partie 1}
Pour faire clignoter la LED, il suffit de faire un XOR sur la bonne addresse du registre P1OUT. Il n'y a pas eu de difficultés majeures sur cette étape.

\lstinputlisting{Code1.c}

\subsection{Partie 2}
Pour la seconde partie il a fallu créer une interruption afin de pouvoir changer l'état des LEDs 1 à 3.\\
Pour ce faire, il faut tout d'abord préparer le bouton en input, ensuite activer la résistance interne, la mettre en "pull-up" et faire en sorte qu'elle détecte correctement les flancs en modifiant le registre approprié.
Il faut ensuite autoriser les interruptions sur les deux entrée des boutons, puis autoriser les interruptions de manières générale.\\
L'étape suivante est d'écrire les fonctions d'interruptions qui vont agir sur les registres. Il faut les précéder de la directive \#pragma afin d'indexer les registres des ports 1 et 2.
Il ne faut pas oublier de remettre à zéro le bit du registre contenant le flag d'interruption sinon il n'est pas possible de continuer l'exécution du programme principal.

La principale difficulté que j'ai rencontré dans cette étape de la conception était l'oubli de la résistance interne rendant le bouton inutilisable pendant quelques secondes après chaque appui, le temps que la tension se dissipe.
Un autre difficulté fût d'implémenter un anti-rebond pour ne pas prendre un compte les appuis "parasites" sur le bouton.

\lstinputlisting{Code2.c}

\subsection{Partie 3}
Cette partie de la conception s'est révélée relativement similaire à la première. La principale différence se situtant dans la fonction d'interrution.

\lstinputlisting{Code3.c}

\subsection{Partie 4}
Il s'agit du même code que pour la partie 2 mais avec un seul bouton.

\lstinputlisting{Code4.c}


\part{Conclusion}
\section{Conclusion}
En conclusion, il est pratique de pouvoir utiliser les interruptions pour effectuer des actions simples sans perturber outre mesure la boucle principale.
Il faut toutefois rester prudent et s'assurer que le code contenu dans ces fonctions est simple et rapide à exécuter.


\end{document}
